\chapter{Ограничение сложности конечной геометрии}

После отрицательных проверок носителя, треугольника редукций и граничного
источника остаётся класс действительно новой конечной геометрии. Чтобы его
исследование не превратилось в неограниченное добавление прямых слагаемых и
стрелок, заранее фиксируется конечная область перебора.

\section{Граница, заданная литературой}

Конечные спектральные тройки классифицируются диаграммами Краевского; см.
\path{arXiv:hep-th/9701081}. Известны также переборы геометрий KO-размерности
6 с четырьмя и шестью слагаемыми; см. \path{arXiv:hep-th/0610040} и
\path{arXiv:0901.3214}. Поэтому новизна проекта может состоять лишь в
конкретном сочетании ограничений, а не в самом переборе диаграмм.

\section{Зафиксированная область перебора}

\begin{definition}[Конечная область перебора версии V]
Первый цикл ограничивается обычными нескрученными вещественными конечными
спектральными тройками:
\begin{enumerate}
 \item KO-размерность 6, вещественная структура, градуировка, условия
 нулевого и первого порядка;
 \item не более пяти простых слагаемых \(M_n(\mathbb K)\), где
 \(\mathbb K=\mathbb R,\mathbb C,\mathbb H\) и \(n\le4\);
 \item не более пяти вершин частичного сектора, шести неориентированных
 рёбер и кратности ребра не более трёх;
 \item связный двудольный остов частичного сектора;
 \item размерность частичного гильбертова пространства не более \(64\);
 \item цикломатическое число не более единицы;
 \item скалярный коммутант и отсутствие независимых центральных весов следа.
\end{enumerate}
\end{definition}

Двенадцать типов простых алгебр и не более пяти неупорядоченных слагаемых
дают не более
\begin{equation}
 \sum_{s=1}^{5}\binom{12+s-1}{s}=6187
\end{equation}
наборов алгебр до отсечения по размерности и представлению.

\section{Перебор графов}

Числа связных простых двудольных неизоморфных графов равны
\begin{equation}
 N_3=1,\qquad N_4=3,\qquad N_5=5.
\end{equation}
При четырёх вершинах получаются остовы \(K_{1,3}\), \(P_4\) и \(C_4\).
Деревья не содержат циклов. Квадрат уже проверялся в двух вариантах:
условие первого порядка сохраняет фазу, но пассивные кратности поколений
остаются произвольными, а аффинное поднятие \(C_4\) смешивает канонические
синглет и триплет.

Для пяти вершин реестр последовательностей степеней и цикломатических чисел
имеет вид
\begin{center}
\small
\begin{tabular}{ccc}
\toprule
Последовательность степеней & Число рёбер & Цикломатическое число\\
\midrule
\((4,1,1,1,1)\) & 4 & 0\\
\((3,2,1,1,1)\) & 4 & 0\\
\((2,2,2,1,1)\) & 4 & 0\\
\((3,2,2,2,1)\) & 5 & 1\\
\((3,3,2,2,2)\) & 6 & 2\\
\bottomrule
\end{tabular}
\end{center}

Следовательно, единственным новым остовом с одним циклом является
\begin{equation}
 K_{2,3}\setminus e,
\end{equation}
то есть квадрат с одной висячей вершиной.

\begin{proposition}[Предварительный выбор графа]
В заданной области \(K_{2,3}\setminus e\) является единственным минимальным
пятивершинным остовом с одним циклом. Это лишь комбинаторный результат.
Вопрос о том, ограничивает ли висячая вершина какой-либо циклический блок,
решается отдельно условием первого порядка.
\end{proposition}

\section{Алгебра предполагаемого селектора}

Для устранения пассивной матричной свободы трёх поколений минимальными
неприводимыми ассоциативными действиями являются \(M_3(\mathbb R)\) на
вещественном триплете и \(M_3(\mathbb C)\) на комплексном триплете.
Вещественный вариант рассматривается первым, поскольку согласуется с уже
имеющимися тетраэдрическим носителем и модой Майораны и имеет меньшую
размерность алгебры.

\begin{nogocriterion}[Запрет перехода от графа к селектору]
Висячая вершина считается селектором лишь тогда, когда уравнения первого
порядка вынуждают уже существующее ребро цикла принадлежать скалярному
коммутанту. Одного присоединения модуля \(M_3(\mathbb R)\) недостаточно.
\end{nogocriterion}

\section{Вывод}

\begin{protocol}
\begin{enumerate}
 \item конечная область перебора: \textbf{задана};
 \item графы до пяти вершин: \textbf{перечислены};
 \item единственный предварительный одноцикловый остов: \textbf{найден};
 \item действие селектора и существование конечной геометрии:
 \textbf{не установлены};
 \item физическое замыкание: \textbf{не достигнуто}.
\end{enumerate}
\end{protocol}

Следующая проверка --- \path{version5_real_selector_leaf_ko6_gate}.
Результат машинного перебора сохранён в
\path{s2t_v5_finite_geometry_complexity_bound_gate_results.json}.